\section{결론 및 제언}

% --------------------------------------------------------------------
% 이 파일의 역할
% --------------------------------------------------------------------
% 5-Conclusions.tex에는 연구 전체의 결론과 제언을 씁니다.
% 연구 목적과 연구 문제에 대한 최종 답을 제시하고,
% 연구 결과가 학문적, 교육적, 실천적으로 어떤 의미를 갖는지 일반화합니다.
% 또한 연구의 의의, 제한점, 후속 연구나 교육 현장을 위한 제언을 정리합니다.
%
% --------------------------------------------------------------------
% 학위논문 결론 장 작성 안내
% --------------------------------------------------------------------
% 결론 장은 연구 전체를 닫는 부분입니다. 단순히 연구 결과를 다시 요약하는 장이 아니라,
% 연구 목적에 대해 최종적으로 어떤 답을 얻었는지 밝히는 장입니다.
%
% 가장 중요한 원칙
% 1. 결론에는 연구 목적에 대한 답이 있어야 합니다.
%    서론에서 제시한 연구 목적과 연구 문제가 결론에서 다시 회수되어야 합니다.
%    즉, "이 연구는 무엇을 밝히고자 했고, 그 결과 무엇을 밝혔다"가 드러나야 합니다.
%
% 2. 결론은 결과를 일반화해야 합니다.
%    결과 장이 구체적인 수치, 표, 그림, 사례를 제시하는 곳이라면,
%    결론 장은 그 결과가 학문적, 교육적, 실천적으로 무엇을 의미하는지
%    더 넓은 수준의 문장으로 정리하는 곳입니다.
%    다만 연구 대상과 방법의 범위를 넘어서는 과도한 일반화는 피해야 합니다.
%
% 3. 결론은 연구 문제의 순서와 대응되게 쓰는 것이 좋습니다.
%    연구 문제 1에 대한 결론, 연구 문제 2에 대한 결론처럼 정리하면
%    독자가 연구 목적, 결과, 결론의 연결을 쉽게 확인할 수 있습니다.
%
% 작성 순서 예시
% 1. 연구 목적과 연구 절차의 간단한 재진술
% 2. 주요 결론
% 3. 연구의 의의
% 4. 연구의 제한점
% 5. 후속 연구 또는 현장 적용을 위한 제언

결론 및 제언 장은 연구 전체를 닫으면서 과학교육 연구 목적에 대한 최종 답을 제시하는 곳이다.
결과 장의 표와 수치를 그대로 반복하기보다, 그 결과가 과학 개념 학습, 탐구 활동, 교수·학습 설계, 평가, 교사 교육 또는 교육과정 실행 측면에서 무엇을 의미하는지 더 넓은 문장으로 정리한다.
다만 연구 대상과 방법의 범위를 넘어서는 과도한 일반화는 피해야 한다.

이 템플릿에서는 결론 및 제언을 \verb|sub/5-Conclusions.tex|에 작성한다.
결론 장도 다른 장과 마찬가지로 \verb|\subsection{}|으로 연구 요약, 주요 결론, 의의, 제한점, 제언을 나누어 쓴다.
필요하면 앞 장의 표나 그림을 \verb|\ref{}|로 다시 참조할 수 있지만, 결론 장에서는 번호와 수치를 길게 반복하기보다 핵심 의미를 서술하는 편이 좋다.

\subsection{연구 요약}

% 이 절에서는 연구 목적, 연구 대상, 연구 방법을 매우 간단히 다시 제시합니다.
% 서론이나 연구 방법 장의 내용을 길게 반복하지 말고, 결론을 이해하는 데 필요한
% 최소한의 맥락만 정리합니다.

본 연구의 목적은 연구자가 설정한 과학교육 연구 문제에 답하고, 그 결과를 바탕으로 해당 교수·학습 현상에 대한 이해를 확장하는 데 있다.
이를 위해 연구 목적에 적합한 수업 자료, 검사 결과, 설문 응답, 면담 자료 또는 관찰 기록을 수집하고 분석하였으며, 분석 결과를 연구 문제의 순서에 따라 제시하였다.

\subsection{주요 결론}

% 이 절에는 연구 목적과 연구 문제에 대한 최종 답을 씁니다.
% 결과 장의 표와 수치를 그대로 반복하기보다, 결과가 말해 주는 핵심 의미를
% 일반화된 문장으로 정리합니다.
%
% 문장 예시
% - 첫째, 본 연구의 결과는 ...임을 보여준다.
% - 둘째, ...는 ...에 긍정적인 영향을 미칠 수 있음을 시사한다.
% - 셋째, ...의 차이는 ...라는 점에서 해석될 수 있다.
%
% 좋지 않은 방식
% - 표 1에서 평균은 82.15점이었다.
% - 실험 집단은 13.75점 증가하였다.
%
% 좋은 방식
% - 이러한 결과는 해당 교수·학습 방법이 학습자의 이해 향상에 기여할 가능성을 보여준다.
% - 집단 간 변화 폭의 차이는 수업 처치가 학습 경험의 질을 바꾸는 요인으로 작용했을 가능성을 시사한다.

첫째, 연구 결과는 연구 목적에서 제기한 핵심 질문에 대한 답을 제공한다.
구체적인 분석 결과를 종합하면, 본 연구에서 다룬 과학 개념, 탐구 활동, 교수·학습 처치 또는 학습자 반응은 연구 맥락 안에서 의미 있는 양상을 보였다.

둘째, 이러한 결과는 개별 사례나 수치에 머무르지 않고, 유사한 과학 수업 맥락이나 과학교육 연구 맥락을 이해하는 데 활용될 수 있다.
따라서 결론에서는 결과 장의 세부 값을 반복하기보다, 그 결과가 더 넓은 맥락에서 어떤 의미를 갖는지 일반화하여 서술한다.

\subsection{연구의 의의}

% 연구의 의의는 이 연구가 기존 연구, 이론, 방법, 교육 현장, 정책, 실천에
% 어떤 기여를 하는지 설명하는 부분입니다.
% "새롭다", "중요하다"라고만 쓰지 말고, 무엇이 왜 새롭고 중요한지 구체적으로 씁니다.

본 연구는 연구 주제와 관련된 과학교육 논의를 보완하고, 후속 연구와 과학 수업 현장 적용을 위한 기초 자료를 제공한다는 점에서 의의가 있다.
또한 연구 결과를 바탕으로 학습자의 이해, 탐구 수행, 수업 설계 또는 평가에 대한 이해를 구체화하고, 실제 적용 가능성을 검토할 수 있는 근거를 제시한다.

\subsection{연구의 제한점}

% 제한점은 연구의 가치를 낮추는 부분이 아니라, 연구 결과를 정확하게 이해하기 위한 범위를 밝히는 부분입니다.
% 표본 규모, 연구 대상, 연구 기간, 측정 도구, 자료 수집 맥락, 분석 방법의 한계를 정직하게 씁니다.
% 제한점을 쓴 뒤에는 그 제한점이 결과 해석에 어떤 영향을 줄 수 있는지도 설명합니다.

본 연구는 연구 대상, 자료 수집 기간, 분석 방법의 범위 안에서 해석될 필요가 있다.
따라서 연구 결과를 모든 상황에 그대로 적용하기보다는, 본 연구가 수행된 맥락과 조건을 고려하여 이해해야 한다.

\subsection{제언}

% 제언은 결론에서 도출된 내용을 바탕으로 씁니다.
% 연구 결과와 무관한 일반적 주장이나 당위적 문장을 나열하지 않도록 주의합니다.
% 보통 다음 두 가지 방향으로 씁니다.
% 1. 후속 연구를 위한 제언
% 2. 교육 현장, 정책, 프로그램 개발 등을 위한 실천적 제언

후속 연구에서는 본 연구의 제한점을 보완하여 더 다양한 학교급, 학년, 과학 과목, 수업 단원, 학습자 집단에서 연구 결과를 검토할 필요가 있다.
또한 본 연구에서 확인한 결과를 바탕으로 실제 과학 수업, 교수·학습 자료 개발, 평가 도구 개선 또는 교사 연수 맥락에서 적용 가능성을 탐색할 수 있다.
